\begin{recipe}
[%
    preparationtime = {\unit[45]{min}},
    bakingtime = {\unit[30]{min}},
    portion = {\portion{6}},
    source = {\protect\recipesource{https://emmikochteinfach.de/lasagne-rezept/}{Emmi kocht einfach: Lasagne}},
    calory = {560kcal | 32g Protein | 47g KH | 31g Fett},
    price = {2,00€},
    created = {05.07.2026},
    %rating = {1},
]
{Lasagne mit Tofuhack}

\graph{% pictures
    big=pic/hauptgericht/19_lasange
}

\introduction{%
Eine klassische Lasagne mit kräftiger Tomaten-Bolognese, cremiger Béchamelsauce und goldbrauner Parmesankruste. In dieser fleischfreien Variante ersetzt Tofuhack das Hackfleisch; auf Speck wird komplett verzichtet.
}

\ingredients{
    \textbf{Lasagne} & \\
    8--12 & Lasagneplatten, ungekocht\index{Getreide!Pasta}\\
    150 g & Parmesan, frisch gerieben\index{Milchprodukte \& Alternativen!Parmesan}\\
    & Butter für Form und Topping\\
    \textbf{Bolognese} & \\
    400 g & Tofuhack\index{Tofu \& Fleischalternativen!Tofuhack}\\
    400 g & Passierte Tomaten\index{Gemüse!Tomate}\\
    2--3 & Zwiebeln, ca. 150 g\\
    1 & Möhre, ca. 50 g\index{Gemüse!Karotte}\\
    1 Stange & Sellerie, ca. 50 g\index{Gemüse!Sellerie}\\
    200 ml & Vollmilch\index{Milchprodukte \& Alternativen!Milch}\\
    40 g & Butter\\
    1 EL & Pflanzenöl\\
    1 EL & Tomatenmark\\
    1 & Lorbeerblatt\\
    1 Msp. & Muskatnuss\\
    & Salz, schwarzer Pfeffer\\
    \textbf{Béchamel} & \\
    500 ml & Milch\index{Milchprodukte \& Alternativen!Milch}\\
    25 g & Mehl\\
    25 g & Butter\\
    & Salz, schwarzer Pfeffer, Muskatnuss
}

\preparation{%
    \step%
    Backofen auf 200\textcelcius{} Ober-/Unterhitze vorheizen. Parmesan fein reiben und bis zur Verwendung kalt stellen. Zwiebeln und Möhre schälen, Sellerie putzen und alles fein würfeln.

    \step%
    Für die Bolognese Pflanzenöl und 40 g Butter in einem großen Topf erhitzen. Zwiebeln, Möhre und Sellerie darin bei mittlerer Hitze etwa 4 Minuten anschwitzen.

    \step%
    Tofuhack dazugeben und 4--5 Minuten mitbraten. Mit Salz und Pfeffer würzen, dann das Tomatenmark kurz mitrösten.

    \step%
    Mit 200 ml Vollmilch ablöschen, Muskatnuss einrühren und die Flüssigkeit unter gelegentlichem Rühren fast vollständig einkochen lassen.

    \step%
    Passierte Tomaten und Lorbeerblatt unterrühren. Die Sauce mit halb aufgelegtem Deckel mindestens 20 Minuten köcheln lassen. Wird sie zu dick, etwas Milch oder Wasser ergänzen. Lorbeerblatt entfernen und abschmecken.

    \step%
    Für die Béchamel 25 g Butter in einem zweiten Topf schmelzen. Mehl mit einem Schneebesen einrühren und etwa 2 Minuten anschwitzen, ohne es bräunen zu lassen.

    \step%
    Nach und nach 500 ml Milch einrühren. Kurz aufkochen und unter regelmäßigem Rühren köcheln lassen, bis die Sauce cremig eindickt. Kräftig mit Salz, Pfeffer und Muskatnuss abschmecken.

    \step%
    Eine etwa 26~x~19 cm große Auflaufform buttern und 1--2 EL Béchamel auf dem Boden verstreichen. Mit Lasagneplatten auslegen; bei Bedarf passend brechen.

    \step%
    Darauf nacheinander Bolognese, 2--3 EL Béchamel und etwas Parmesan verteilen. Die Schichten wiederholen, bis die Zutaten aufgebraucht sind. Mit reichlich Béchamel und dem restlichen Parmesan abschließen und ein paar Butterflöckchen daraufgeben.

    \step%
    Auf der zweiten Schiene von unten etwa 30 Minuten backen, bis die Oberfläche goldbraun ist. Vor dem Anschneiden 10 Minuten ruhen lassen.
}

\hint{
    Die Ruhezeit nach dem Backen lohnt sich: Die Schichten setzen sich, die Lasagne bleibt beim Portionieren stabiler und schmeckt trotzdem noch richtig heiß.
}

\suggestion[Bolognese]{%
    Wenn genug Zeit da ist, die Tomatensauce länger als 20 Minuten sanft köcheln lassen. Das macht sie kräftiger und runder; bei Bedarf zwischendurch etwas Wasser oder Milch nachgießen.
}

\end{recipe}
